\section*{Chapter \ref{ch:g4:animal-reproduction} --- How Animals Reproduce}

\begin{solution}{exo:g4:animal-reproduction:1}
The joining of one \omterm{def:g4:animal-reproduction:sexual}{sperm cell} with one \omterm{def:g4:animal-reproduction:sexual}{egg cell}, which starts a
new young \omterm{def:g1:animals-around-us:animal}{animal}. The \omterm{def:g4:animal-reproduction:sexual}{egg cell} comes from the \omterm{def:g4:animal-reproduction:sexual}{female}, the \omterm{def:g4:animal-reproduction:sexual}{sperm
cell} from the \omterm{def:g4:animal-reproduction:sexual}{male}.
\end{solution}

\begin{solution}{exo:g4:animal-reproduction:2}
\omterm{def:g4:animal-reproduction:ovivivi}{Oviparous}: the \omterm{def:g4:animal-reproduction:sexual}{female} lays \omterm{def:g2:animals-grow-up:born}{eggs} and the young finishes growing
in the \omterm{def:g2:animals-grow-up:born}{egg}, outside her body --- the hen, the frog (also trout,
lizards, snails, \omterm{prop:g3:sorting-living-things:groups}{insects}). \omterm{def:g4:animal-reproduction:ovivivi}{Viviparous}: the young grows inside
the mother, fed by her, and is \omterm{def:g2:animals-grow-up:born}{born alive} --- the mare, the
whale (nearly all \omterm{prop:g3:sorting-living-things:groups}{mammals}).
\end{solution}

\begin{solution}{exo:g4:animal-reproduction:3}
The \omterm{def:g4:animal-reproduction:sexual}{males} sang to call the \omterm{def:g4:animal-reproduction:sexual}{females} to the pond; each pair
released its cells into the water, where sperm met \omterm{def:g2:animals-grow-up:born}{eggs} ---
\omterm{def:g4:animal-reproduction:sexual}{fertilization}; the fertilized \omterm{def:g2:animals-grow-up:born}{eggs} in their jelly grew into the
June \omterm{ex:g2:animals-grow-up:frog}{tadpoles}.
\end{solution}

\begin{solution}{exo:g4:animal-reproduction:4}
The trout's happens in the water, outside the \omterm{def:g4:animal-reproduction:sexual}{female} --- \omterm{def:g2:animals-grow-up:born}{eggs}
and sperm are shed together. The hen's happens inside her body,
during \omterm{prop:g4:animal-reproduction:where}{mating}. On land the naked cells would dry out; in water
they can safely meet outside.
\end{solution}

\begin{solution}{exo:g4:animal-reproduction:5}
The chick grows in a shelled \omterm{def:g2:animals-grow-up:born}{egg} outside the hen, living off
the yolk packed with it. The foal grows eleven months inside
the mare, fed continuously by her body, and is \omterm{def:g2:animals-grow-up:born}{born alive}.
\end{solution}

\begin{solution}{exo:g4:animal-reproduction:6}
The more young a species produces, the less care each gets;
the fewer, the more care. Extremes: the codfish's millions of
uncared-for \omterm{def:g2:animals-grow-up:born}{eggs}, and the elephant's single calf guarded for a
decade.
\end{solution}

\begin{solution}{exo:g4:animal-reproduction:7}
Because success only requires that, on average, enough young
survive to replace the parents. Out of millions of \omterm{def:g2:animals-grow-up:born}{eggs}, a few
always slip through the sea's dangers --- the huge number
\emph{is} the protection.
\end{solution}

\begin{solution}{exo:g4:animal-reproduction:8}
Few young (five is few compared to thousands), grown inside the
mother, nursed and guarded: the cat sits clearly on the
few-and-cared-for side --- like the \omterm{prop:g3:sorting-living-things:groups}{mammals} generally.
\end{solution}

\begin{solution}{exo:g4:animal-reproduction:9}
That it receives something from \emph{each} parent --- the
young is not a copy of its mother alone. Both parents
contribute to what the foal becomes.
\end{solution}

\begin{solution}{exo:g4:animal-reproduction:10}
A small clutch --- some dozens or hundreds of \omterm{def:g2:animals-grow-up:born}{eggs}, not
millions. Care and numbers trade off: a father who guards,
fans and defends can bring a small clutch through dangers that
would demand millions of unguarded \omterm{def:g2:animals-grow-up:born}{eggs}.
\end{solution}

\begin{solution}{exo:g4:animal-reproduction:11}
\omterm{def:g4:animal-reproduction:ovivivi}{Oviparous} --- and uncared for: the buried clutch is left to
the sand's warmth, and the hatchlings dash to the sea alone
(their large numbers are their strategy). She must lay on land
because reptile \omterm{def:g2:animals-grow-up:born}{eggs} are shelled and breathe air --- laid in
the sea they would drown. Her class's features, not her
habits, set the rule.
\end{solution}
